\documentclass[../main.tex]{subfiles}

\setcounter{section}{16}

\begin{document}

\begin{theorem}[Division Algorithm for $\poly{F}$]
  Let $F$ be a field and $f(x) \ g(x) \in \poly{F}$ where $g(x) \neq 0$.
  Then there are unique $q(x) \ r(x) \in \poly{F}$ such that
  $f(x) = g(x)q(x) + r(x)$ and either $r(x) = 0$ or $\deg r(x) < \deg g(x)$.
\end{theorem}
\begin{proof}
  % If $\deg f(x) < \deg g(x)$, then $q(x) = 0$ and $r(x) = f(x)$.
  % Now suppose $\deg f(x) = m$ and $\deg g(x) = n$ where $m \ge n$.
  % There exist $q$ such that $f_m = g_n \times q$ where
  % $f_m \ g_n$ are coefficients of $f(x) \ g(x)$ respectively.
  % We claim $q$ is the coefficient of degree $m - n$ of $q(x)$,
  % we continue this progress on $f(x) - (g(x) \times q x^{m - n})$ 
  % and find other coefficients,
  % note that the degree is smaller than $m$, so the algorithm will terminate at some point.

  If $\deg f(x) < \deg g(x)$, then $q(x) = 0$ and $r(x) = f(x)$, 
  otherwise, we induction/recursion on the degree of $f(x)$:
  Let $m = \deg f(x)$ and $n = \deg g(x)$,
  \begin{itemize}
    \item Base: For any $g(x)$, if $0 \ge n$ which implies $n = 0$, 
      then there is $q$ such that $f(x) = g(x) \times q + 0$ and $0 = 0$.
    \item Induction: For any $f(x)$ with degree belows $m$, and any $g(x)$,
      there are $q(x)$ and $r(x)$ such that $f(x) = g(x)q(x) + r(x)$,
      then there are $q$ such that $f_m = g_n \times q$, also, by induction hypothesis,
      we know there are $q^\prime(x)$ and $r(x)$ 
      such that $f(x) - g(x)qx^{m - n} = g(x)q^\prime(x) + r(x)$,
      which is in fact $f(x) = g(x)q^\prime(x) + r(x) + g(x) qx^{m - n} = g(x) (q^\prime(x) + qx^{m - n}) + r(x)$.
  \end{itemize}

  The result is unique by comparing the degree.
\end{proof}

\begin{corollary}[Remainder Theorem]
  Let $F$ be a field, $a \in F$ and $f(x) \in \poly{F}$. Then $f(a)$
  is the remainder in the division of $f(x)$ by $x - a$.
\end{corollary}
\begin{proof}
  Induction on the degree of $f(x)$ with hypothesis:
  Let $n = \deg \poly{F}$, $m \in F$, then $f(a) + ma^{n + 1}$
  is the remainder in the division of $f(x) + max^n$ by $x - a$.
  \begin{itemize}
    \item Base: Let $n = 0$, $m \in F$, then the remainder
          in the divison of $f(x) + max^0 = f_0 + ma$ by $x - a$
          is obviously $f_0 + ma = f(a) + ma^{0 + 1} = f(a) + ma$.
    \item Induction: Let $n = \deg f(a)$, $m \in F$, the remainder in 
          the division of $f(x) + max^n$ by $x - a$ is
          the divison of $f(x) + max^n - (f_n + ma)x^{n - 1}(x - a) = f^\prime(x) + (f_n + ma)x^{n - 1}a$ by $x - a$
          where $f^\prime(x) = f(x) - f_n x^n$.
          By induction hypothesis, we know the remainder in the divison of $f^\prime(x) + (f_n + ma)x^{n - 1}$
          by $x - a$ is $f^\prime(a) + (f_n + ma)a^{n - 1 + 1} = f^\prime(a) + (f_n + ma)a^n = f(a) + ma^{n + 1}$.
  \end{itemize}

  Then, choose $m = 0$, we get $f(a)$ is the remainder is the divison of $f(a)$ by $x - a$.
\end{proof}

\end{document}